\begin{mybox}
\section{Python Testing with pytest}
\label{sec:org2f60ca8}
\tcblower{}
\url{https://learning.oreilly.com/library/view/python-testing-with/9781680509427/f\_0012.xhtml}
\end{mybox}\begin{myitems}
\begin{itemize}
\item 7.5 Writing a Test Strategy
\item 8.1 Understanding pytest Configuration Files
\item 8.2 Saving Settings and Flags in \texttt{pytest.ini}
\item 8.3 Using \texttt{tox.ini}, \texttt{pyproject.toml}, or \texttt{setup.cfg} in place of \texttt{pytest.init}
\item 8.4 Determining a Root Directory and Config File
\item 8.5 Sharing Local Fixtures and Hook Functions with \texttt{conftest.py}
\item 8.6 Avoiding Test File Name Collision
\item 9.1 Using \texttt{coverage.py} with \texttt{pytest-cov}
\item 9.2 Generating HTML Reports
\item 9.3 Excluding Code from Coverage
\item 9.4 Running Coverage on Tests
\item 9.5 Running Coverage on Directory
\item 9.6 Running on a Single File
\item 10.1 Isolating the Command-Line interface
\item 10.2 Testing with Typer
\item 10.3 Mocking an Attribute
\item 10.4 Mocking a Class and Methods
\item 10.5 Keeping Mock and Implementation in Sync with \texttt{Autospec}
\item 10.6 Making Sure Functions Are Called Correctly
\item 10.7 Creating Error Conditions
\item 10.8 Testing at Multiple Layers to Avoid Mocking
\item 10.9 Using Plugins to Assist Mocking
\end{itemize}

\end{myitems}\begin{mybox}
\section{Docker in Action}
\label{sec:orgf9c5535}
\tcblower{}
\url{https://learning.oreilly.com/library/view/docker-in-action/9781617294761/kindle\_split\_013.html}
\end{mybox}\begin{myitems}
\begin{itemize}
\item 1.1 What is Docker
\begin{itemize}
\item 1.1.1 \texttt{"Hello, World"}
\item 1.1.2 Containers
\item 1.1.3 Containers are not virtualization
\item 1.1.4 Running software in containers for isolation
\item 1.1.5 Shipping containers
\end{itemize}
\item 1.2 What problems does Docker solve?
\begin{itemize}
\item 1.2.1 Getting organized
\item 1.2.2 Improving portability
\item 1.2.3 Protecting your computer
\end{itemize}
\item 1.3 Why is Docker important
\item 1.4 Where and when to use Docker
\item 1.5 Docker in the larger ecosystem
\item 1.6 Getting help with the Docker command line
\end{itemize}

\end{myitems}\begin{mybox}
\section{Introduction to HTML}
\label{sec:org5ceddf3}
\tcblower{}
\url{https://developer.mozilla.org/en-US/docs/Learn/HTML/Introduction\_to\_HTML}
\end{mybox}\begin{myitems}
\begin{itemize}
\item 1.4 Creating hyperlinks
\begin{itemize}
\item 1.4.1 What is a hyperlink?
\item 1.4.2 Anatomy of a link
\item 1.4.3 A quick primer on URLs and paths
\item 1.4.4 Link best practices
\item 1.4.5 E-mail links
\end{itemize}
\item 1.5 Advanced text formatting
\begin{itemize}
\item 1.5.1 Description lists
\item 1.5.2 Quotations
\item 1.5.3 Abbreviations
\item 1.5.4 Marking up contact details
\item 1.5.5 Superscript and subscript
\item 1.5.6 Representing computer code
\item 1.5.7 Marking up times and dates
\end{itemize}
\item 1.6 Document and website structure
\begin{itemize}
\item 1.6.1 Basic sections of a document
\item 1.6.2 HTML for structuring content
\item 1.6.3 HTML layout elements in more detail
\item 1.6.4 Planning a simple website
\end{itemize}
\item 1.7 Debugging HTML
\item 1.8 Assessment: Marking up a letter
\item 1.9 Assessment: Structuring a page of content
\item 2.1 Multimedia and embedding overview
\item 2.2 Images in HTML
\end{itemize}

\end{myitems}\begin{mybox}
\section{AWS Essential Training for Developers}
\label{sec:orgdc8776e}
\tcblower{}
\url{https://www.linkedin.com/learning/aws-essential-training-for-developers}
\end{mybox}\begin{myitems}
\begin{itemize}
\item 9.1 What is DevOps?
\item 9.2 CI/CD with AWS
\item 9.3 Infrastructureas code with AWS
\item 9.4 Monitoring with CloudWatch
\item 10.1 AWS Shield and firewalls with WAF
\item 10.2 Inspector, GuardDuty, and Macie
\item 10.3 CloudTrail and Security Hub
\end{itemize}

\end{myitems}

除此之外，还有在leetcode(\url{https://leetcode.com})刷题。
